\documentclass{article}

\usepackage{amsmath,amssymb}
\usepackage{xurl}
\usepackage[hidelinks]{hyperref}
\setlength{\emergencystretch}{2em}

% Shared commands for notes in docs/paper-gaps/.

\DeclareUnicodeCharacter{2080}{\ifmmode{}_{0}\else\textsubscript{0}\fi}
\DeclareUnicodeCharacter{2081}{\ifmmode{}_{1}\else\textsubscript{1}\fi}
\DeclareUnicodeCharacter{2082}{\ifmmode{}_{2}\else\textsubscript{2}\fi}
\DeclareUnicodeCharacter{2099}{\ifmmode{}_{n}\else\textsubscript{n}\fi}

\newcommand{\Lean}{\textsc{Lean}}
\ExplSyntaxOn
\NewDocumentCommand{\leanid}{m}{
  \ifmmode
    \textsf{\detokenize{#1}}
  \else
    \group_begin:
    \urlstyle{sf}
    \tl_set:Nn \l_tmpa_tl {#1}
    \tl_replace_all:Nnn \l_tmpa_tl {\_} {_}
    \exp_args:NV \path \l_tmpa_tl
    \group_end:
  \fi
}
\ExplSyntaxOff
\providecommand{\tr}{\operatorname{tr}}
\newcommand{\tnleanIssueBase}{https://github.com/LionSR/TNLean/issues/}
\newcommand{\tnleanPrBase}{https://github.com/LionSR/TNLean/pull/}
\newcommand{\tnleanDocBase}{https://sirui-lu.com/TNLean/docs/}
\newcommand{\ghissue}[1]{\href{\tnleanIssueBase#1}{GitHub issue \##1}}
\newcommand{\ghpr}[1]{\href{\tnleanPrBase#1}{PR \##1}}
\newcommand{\doclink}[1]{\href{\tnleanDocBase#1.html}{\path{#1}}}

% Dirac notation and the identity operator, as in blueprint/src/macros/common.tex.
% \providecommand keeps note-local definitions authoritative.
\usepackage{dsfont}
\providecommand{\ket}[1]{|#1\rangle}
\providecommand{\bra}[1]{\langle #1|}
\providecommand{\braket}[2]{\langle #1 | #2 \rangle}
\providecommand{\ketbra}[2]{|#1\rangle\!\langle #2|}
\providecommand{\Id}{\mathds{1}}
\providecommand{\one}{\mathds{1}}
\providecommand{\C}{\mathbb{C}}
\providecommand{\MN}[1]{M_{#1}(\C)}
\providecommand{\MD}{M_D(\C)}

% External referencing for paper sources.  The build helper writes sanitized
% auxiliary files under build/paper-aux/ and excludes duplicated source labels.
\usepackage{xr}

% Import a generated paper auxiliary file with a caller-supplied label prefix.
% The candidate paths support builds from docs/paper-gaps/, the repository root,
% and build/paper-aux/ itself.
\newcommand{\importpaperaux}[2]{%
  \IfFileExists{../../build/paper-aux/#2.aux}{%
    \externaldocument[#1]{../../build/paper-aux/#2}%
  }{%
    \IfFileExists{build/paper-aux/#2.aux}{%
      \externaldocument[#1]{build/paper-aux/#2}%
    }{%
      \IfFileExists{#2.aux}{%
        \externaldocument[#1]{#2}%
      }{}%
    }%
  }%
}

% General external reference macros
\newcommand{\paperref}[2]{\ref{#1-#2}}
\newcommand{\papereqref}[2]{(\ref{#1-#2})}

% CPSV16 source (1606.00608 / MPDO-22-12-17-2.tex) external document and shortcuts
\importpaperaux{cpsv16-}{1606.00608/MPDO-22-12-17-2-tnlean}
\newcommand{\cpsvref}[1]{\paperref{cpsv16}{#1}}
\newcommand{\cpsveqref}[1]{\papereqref{cpsv16}{#1}}

% Verdict marker: \gapnote{<kind>}{<status>}. Kind is the policy
% classification (clarification, local-correction, scope-restriction,
% unfaithful, false-source, open-gap); status is open, wip (elimination
% underway), resolved, or historical. Machine-read by the site index;
% prints a verdict line.
\newcommand{\gapnote}[2]{\par\noindent\textbf{Verdict:} #1 (#2).\par}


\title{Top Schmidt-rank index in the $n$-positivity threshold of $T_\eta$
(Wolf Chapter 3)}
\author{}
\date{2026-06-24}

\begin{document}
\maketitle

\section{Source statement}

In Chapter~3 of Wolf's lecture notes \cite{Wolf2012Quantum}, the one-parameter
family
\begin{align}
  T_\eta(\rho)=\tr(\rho)\,\mathbf 1-\frac{\rho}{\eta},
  \qquad \eta\in\mathbb R_+,
  \tag{3.11}
\end{align}
on $M_D(\mathbb C)$ (the formal map is written $\tr(\rho)\,\mathbf 1-\eta^{-1}\rho$,
which agrees on $\eta>0$ and is moreover defined at $\eta=0$) is analyzed
through its Choi--Jamiolkowski operator
$\tau_\eta=D^{-1}\mathbf 1-\eta^{-1}\ket\Omega\bra\Omega$.  All positive
eigenvalues of $\tau_\eta$ equal $1/D$, and the single non-positive eigenvalue
is $\nu_-=1/D-1/\eta$ with reduced density $\rho_-=\mathbf 1/D$, so its
Ky--Fan $n$-norm is $\|\rho_-\|_{(n)}=n/D$.  The source concludes:
\begin{align}
  T_\eta\ \text{is $n$-positive}\iff \eta\ge n,
  \qquad n=1,\dots,D,
\end{align}
where the range of $n$ runs up to the dimension $D$, in which case
$n$-positivity is complete positivity.  This proves that every inclusion in
Wolf's chain
\begin{align}
  \mathcal T_{\mathrm{cp}}=\mathcal T_D\subseteq
  \mathcal T_{D-1}\subseteq\cdots\subseteq\mathcal T_2\subseteq\mathcal T_1
  \tag{3.3}
\end{align}
is strict.\footnote{Local source:
\path{Notes/WolfNoteTexSource/ch03_positive_not_completely.tex:181--193},
equation label \path{eq:ch3:3.11}.}

\section{Point at issue}

The formal statement establishes the equivalence
\begin{align}
  T_\eta\ \text{is $k$-positive}\iff \eta\ge k
\end{align}
for Schmidt-rank index $1\le k<D$, omitting the top index $k=D$.  The
restriction is inherited, not introduced here: the equivalence is derived from
the maximal-overlap principle (Wolf Lemma~3.1), whose formal version
$\sup_\psi|\braket{\Omega}{\psi}|^2=\|\rho\|_{(k)}$ over Schmidt-rank-$k$
normalized vectors is available only for $1\le k<D$.  That bound in turn rests
on the Ky--Fan maximum principle, stated for orthogonal projections of rank
exactly $k$ with $k<D$.  At the top index $k=D$ the Ky--Fan $D$-norm is the
full trace and the maximizing projection is the identity; the present
formalization of the maximum principle does not cover that boundary case.

Consequently the strictness witness produced here, a $k$-positive map that fails
to be $(k+1)$-positive, is available for $1\le k$ with $k+1<D$.  This covers
every inclusion in the chain~(3.3) except the topmost one,
$\mathcal T_D\subseteq\mathcal T_{D-1}$, which compares complete positivity with
$(D-1)$-positivity and requires the $k=D$ index.

\section{Status}

The boundary case $k=D$ is resolved by a route that bypasses the $k<D$
restriction of the maximal-overlap and Ky--Fan principles entirely.  On
$M_D(\mathbb C)$, $D$-positivity is complete positivity
(\leanid{isCPMap_iff_isNPositiveMap_card}), which is positive semidefiniteness of
the Choi operator $\tau_\eta$.  Its quadratic form on a vector $\psi$ is
$D^{-1}\|\psi\|^2-\eta^{-1}|\braket{\Omega}{\psi}|^2$, and since $\ket\Omega$ is
normalized the Cauchy--Schwarz bound $|\braket{\Omega}{\psi}|^2\le\|\psi\|^2$
(saturated at $\psi=\Omega$) makes that form nonnegative for every $\psi$ exactly
when $\eta\ge D$.  Hence $T_\eta$ is $D$-positive iff $\eta\ge D$, the equivalence
at the top index.

The formal declarations are \leanid{Matrix.isNPositiveMap_tEta_iff} for the
range $1\le k<D$ and \leanid{Matrix.isNPositiveMap_tEta_card_iff} for the top
index $k=D$, together covering the full equivalence~(3.11) over $1\le k\le D$, in
\doclink{TNLean/Channel/NPositivityChainStrict}.  The strictness witness
\leanid{Matrix.exists_isNPositiveMap_not_succ} remains stated for $k+1<D$; the
strictness of the top inclusion $\mathcal T_D\subsetneq\mathcal T_{D-1}$ follows
from the $k=D-1$ and $k=D$ thresholds but is not recorded as a separate witness.

\section{Classification}

\emph{Resolved.} The full source range $1\le k\le D$ of equivalence~(3.11) is now
formalized: \leanid{Matrix.isNPositiveMap_tEta_iff} for $1\le k<D$ and
\leanid{Matrix.isNPositiveMap_tEta_card_iff} for the completely-positive endpoint
$k=D$, the latter by the Choi-operator route detailed in the Status above.  No
hypothesis foreign to the source is used.

\bibliographystyle{alpha}
\bibliography{references}

\end{document}
